\documentclass[10pt,a4paper]{amsart}
\usepackage[margin=19mm]{geometry}
\usepackage{amsmath,amssymb,mathtools}
\usepackage[scr=rsfs,cal=euler]{mathalfa}
\usepackage{xfrac}
\usepackage[unicode,hidelinks,bookmarks=false]{hyperref}
\newcommand{\ie}{i.e.\ }
\newcommand{\cf}{cf.\ }
\newcommand{\resp}{respectively\ }
\newcommand{\Subsection}{\subsection}
\providecommand{\nobreakdash}{\nobreak\hbox{-}}
\setlength{\emergencystretch}{2em}
\multlinegap0pt
\usepackage{amssymb}
\usepackage{tikz}
\usepackage{xcolor}
\newtheorem{lem}{Lemma}
\title{A note on the equivalence of Gromov boundary and metric boundary}
\author{Vasudevarao Allu, Abhishek Pandey}
\date{Source: arXiv:2405.10273v3 (CC BY 4.0)}
\begin{document}
\maketitle

\noindent\emph{Source and licence.} A note on the equivalence of Gromov boundary and metric boundary. Version arXiv:2405.10273v3 (CC BY 4.0), distributed by arXiv under the Creative Commons Attribution 4.0 International licence, http://creativecommons.org/licenses/by/4.0/, as recorded in the arXiv licence metadata for this exact version and shown on the abstract page https://arxiv.org/abs/2405.10273v3. Full licence notice: \texttt{licenses/arxiv-2405.10273-CC-BY-4.0.txt}.\par\smallskip

\noindent\textit{Editorial scope.} Counted unit: the article's lem continuous (source0.tex lines 520--534). The article's complete original proof of it is delivered with it (source0.tex lines 536--574, one \texttt{proof} environment of 39 lines). No mathematics is written, repaired or re-derived: the statement and the proof are the article's own lines, in the article's own notation, and no retained line is substituted or re-flowed.  Retained uncounted as the source-local context the proof needs: lines 487--489.  Nothing was omitted from any retained span, so the package carries no omission record.  No retained line contains a citation, so the wrapper carries no bibliography and no bibliography entry.  No retained line contains a figure environment, an image-inclusion command or drawing code, so no image is delivered and the package declares no asset dependency.  Editorial formatting only, in addition: an AMS \texttt{amsart} preamble that restates the article's own declarations and macros used by the retained lines, namely the environments \texttt{lem} on separate counters, together with the standard packages those lines need.  The retained environments print in this document as Lemma 1 in order of appearance, whereas the article prints the same items with the numbering of its own full document.  The complete native source of the article as distributed in the arXiv e-print for this exact version is bundled unchanged as \texttt{sources/159281/source0.tex}.
\begin{lem}\label{AV-P7-lem2}
	Let $(\Omega,d)$ be a bounded minimally nice space. Suppose $(\Omega,d)$ is quasihyperbolically visible. If $\gamma$ is a quasihyperbolic geodesic ray in $\Omega$, then $\lim\limits_{t\to \infty}\gamma(t)$ exists in $\partial_d\Omega$.
\end{lem}

\begin{lem}\label{continuous}
	Let $(\Omega,d)$ be a bounded minimally nice space. Suppose $(\Omega,d)$ is quasihyperbolically visible. Fix $x_0\in \Omega$.
	\begin{itemize}
		\item[(i)] Let $\gamma_n:[0,\infty)\to (\Omega,k_{\Omega})$ be quasihyperbolic geodesic rays with $\gamma_n(0)=x_0$ for all $n$.
		If $(\gamma_n)$ converges uniformly on compact subsets to a geodesic ray $\gamma:[0,\infty)\to (\Omega,k_{\Omega})$, then
		\[
		\lim_{t\to \infty}\gamma(t)=\lim_{n\to \infty}\lim_{t\to\infty}\gamma_n(t).
		\]
		\item[(ii)] Let $\gamma_n:[0,T_n]\to (\Omega,k_{\Omega})$ be quasihyperbolic geodesics with $\gamma_n(0)=x_0$ for all $n$, where $T_n\to\infty$.
		If $(\gamma_n)$ converges uniformly on compact subsets to a geodesic ray $\gamma:[0,\infty)\to (\Omega,k_{\Omega})$, then
		\[
		\lim_{t\to \infty}\gamma(t)=\lim_{n\to \infty}\gamma_n(T_n).
		\]
	\end{itemize}
\end{lem}

\begin{proof}
Since $\overline{\Omega}^d$ is compact, it suffices to prove the claim along an arbitrary convergent subsequence of the boundary values.
	
	\noindent\textbf{Proof of (i).}
By Lemma \ref{AV-P7-lem2}, for each $n$, let $p_n:=\lim_{t\to\infty}\gamma_n(t)\in\partial_d\Omega$ and let $q:=\lim_{t\to\infty}\gamma(t)\in\partial_d\Omega$. Passing to a subsequence, we may assume $p_n\to p\in\partial_d\Omega$. We show that $p=q$. Assume for a contradiction that $p\neq q$. Since $\gamma$ lands at $q$, for every $j\in \mathbb{N}$, there exists $s_j\to \infty$ such that for every $t\ge s_j$, $d(\gamma(t),q)<1/j$. In particular, $d(\gamma(s_j),q)<1/j$. Since $\gamma_n\to\gamma$ uniformly on $[0,s_j]$, there exists $N(j)$ such that
	$d(\gamma_n(s_j),\gamma(s_j))<1/j$ for all $n\ge N(j)$, and hence
	\[
	d(\gamma_n(s_j),q)<2/j\qquad\text{for all }n\ge N(j).
	\]
Now choose an increasing sequence of indices $n_j\ge N(j)$ such that $p_{n_j}\stackrel{d}{\to} p$
and $d(p_{n_j},p)<1/j$ for each $j$. For each $j$, since $\gamma_{n_j}(t)\to p_{n_j}$ as $t\to\infty$, pick $t_j>s_j$ such that
$d(\gamma_{n_j}(t_j),p_{n_j})<1/j$. Then
\[
d(\gamma_{n_j}(t_j),p)\le d(\gamma_{n_j}(t_j),p_{n_j})+d(p_{n_j},p)<2/j.
\]
Set $y_j:=\gamma_{n_j}(s_j)$ and $x_j:=\gamma_{n_j}(t_j)$. Then $y_j\to q$ and $x_j\to p$. By quasihyperbolic visibility, there exists a compact set
$K\subset\Omega$ such that every quasihyperbolic geodesic joining $x_j$ to $y_j$ intersects $K$.
In particular, the geodesic segment $\gamma_{n_j}|_{[s_j,t_j]}$ intersects $K$, so we may choose
$c_j\in[s_j,t_j]$ with $\gamma_{n_j}(c_j)\in K$.
Since $\gamma_{n_j}$ is parametrized by $k_\Omega$-arclength and $\gamma_{n_j}(0)=x_0$,
\[
c_j=k_\Omega(x_0,\gamma_{n_j}(c_j))\le \max_{x\in K}k_\Omega(x_0,x)=:M<\infty.
\]
But $c_j\ge s_j\to\infty$, a contradiction. Hence $p=q$, proving (i).
 
	\noindent\textbf{Proof of (ii).}
	Let $x_n:=\gamma_n(T_n)\in\Omega$. Since $\overline{\Omega}^d$ is compact, passing to a subsequence, we may assume that $x_n\to p\in\overline{\Omega}^d$.
	Since $T_n=k_\Omega(x_0,x_n)\to\infty$ and $k_\Omega$ induces the original topology on $\Omega$, we have $p\in\partial_d\Omega$.
	Let $q:=\lim_{t\to\infty}\gamma(t)\in\partial_d\Omega$. Our aim is to show that $p=q$. If not, choose $s_j\to\infty$ such that $d(\gamma(s_j),q)<1/j$ for each $j$.
		Since $\gamma_n\to\gamma$ uniformly on $[0,s_j]$, for each $j$ there exists $N(j)$ such that
		$d(\gamma_n(s_j),q)<2/j$ for all $n\ge N(j)$.
		Also, since $x_n\to p$, for each $j$ there exists $N'(j)$ such that $d(x_n,p)<1/j$ for all $n\ge N'(j)$. Pick an increasing sequence $n_j\ge \max\{N(j),N'(j)\}$ such that $T_{n_j}>s_j$ (this is possible since $T_n\to \infty$). Set $y_j:=\gamma_{n_j}(s_j)$ and $x_j:=\gamma_{n_j}(T_{n_j})$. Then $y_j\to q$ and $x_j\to p$. By quasihyperbolic visibility, there exists a compact set $K\subset\Omega$ such that every quasihyperbolic geodesic joining $x_j$ to $y_j$ intersects $K$.
	In particular, the segment $\gamma_{n_j}|_{[s_j,T_{n_j}]}$ intersects $K$, so choose
	$c_j\in[s_j,T_{n_j}]$ with $\gamma_{n_j}(c_j)\in K$. As above,
	\[
	c_j=k_\Omega(x_0,\gamma_{n_j}(c_j))\le \max_{x\in K}k_\Omega(x_0,x)=M<\infty,
	\]
	but $c_j\ge s_j\to\infty$, a contradiction. Hence $p=q$ and (ii) follows.
\end{proof}

\end{document}
